%! Author = ben
%! Date = 27.11.2023

% Preamble
\documentclass[../main.tex]{subfiles}

% Packages

% Document
\begin{document}
    \chapter{Ergebnisse}


    \section{Einleitung}
    Bei der ersten Nutzung der App wird der Benutzer gebeten, seine Interessen einzutragen.
    Hierdurch ist es möglich, die Inhalte der App auf den Benutzer individuell anzupassen.
    Anschließend kann der Benutzer den aktuellen Stand seines Lebensstils eintragen.
    Hierfür gibt es verschiedene umweltschädliche Aktionen, wie beispielsweise das Autofahren oder das Essen von Fleisch.
    Der Benutzer gibt an, wie oft er diese Aktionen in der Woche durchführt und
    bekommt anschließend ein wöchentliches Ziel vorgeschlagen, welches er erreichen muss.
    Nach Abschluss der Einleitung erhält der Benutzer jeden Tag eine Benachrichtigung mit der Aufforderung, sein heutiges Verhalten einzutragen.
    (Siehe\ \ref{sec:tracker})


    \section{Videos}\label{sec:videos}
    \begin{wrapfigure}{r}{0.4\textwidth}
        \centering{
            \includegraphics[width=0.39\textwidth]{videos.png}
        }
        \caption{Video Feed}\label{fig:videos}
    \end{wrapfigure}
    SaveWorld bietet eine Vielzahl an Videos, die praktische Tipps für den Alltag vermitteln.
    Diese Videos sind in sieben Kategorien unterteilt: Mobilität, Kochen, Ernährung, Freizeit, Haushalt, Einkauf, Urlaub und Kleidung.
    Der Benutzer hat mehrere Möglichkeiten, um auf diese Videos zuzugreifen:
    \begin{enumerate}
        \item \textbf{Kanäle}: Jede Kategorie hat einen eigenen Kanal, in dem alle Videos dieser Kategorie angezeigt werden.
        So können Nutzer gezielt nach Videos suchen, die sie interessieren.
        \item \textbf{Feed}: Der Feed zeigt dem Benutzer Videos an, die auf seine Interessen zugeschnitten sind.
        Das Feed besteht zu 60\% aus Videos, die den Interessen des Benutzers entsprechen und zu 40\% aus anderen Videos.
    \end{enumerate}
    Jedes angesehene Video kann bewertet werden, sodass die Verbesserung künftiger Videos ermöglicht wird.
    Außerdem ist das Kommentieren und Teilen von Videos möglich.
    Damit der Benutzer sich unabhängig und fachlich fundiert informieren kann, sind unter jedem Video Quellen angegeben.
    Hierbei handelt es sich um differenzierte Quellen um sicherzustellen, dass die Informationen korrekt sind.

    \newpage
    \section{Tracker}\label{sec:tracker}

\end{document}
